A body that learns from experience must do more than persist. Its completed interactions must change a state that can influence what happens next, and the resulting action must be useful under a stated task. In artificial life, that state may share a carrier with the material that moves, consumes energy, deforms and dies. The central question is therefore how acquired content contributes to a continuing physical future, and under which sensory and physiological conditions that contribution remains useful. We address this question by separating acquisition, subsequent use, revision after changed consequences and persistence through material renewal.

Continuous cellular automata provide a setting in which these distinctions can be studied directly. Lenia and its extensions produce localized, moving and interacting structures through local update rules\cite{chan,expandedlenia,plantec}. Search over update parameters and initial states can discover robust sensorimotor behaviour\cite{hamon,asal}. Neural cellular automata likewise support growth, regeneration and developmental organization\cite{growingca,developmental,bricks}. Such capacities establish rich physical and computational possibilities, but they do not identify where learning takes place. A rule chosen by an outer optimizer, a state acquired during one body's lifetime and a fixed physiological response can all generate history-dependent behaviour. They make different predictions when the body encounters an unannounced change or its acquired state is manipulated.

Runtime plasticity offers a more direct route to individual learning. Local Hebbian and homeostatic changes have been used in adaptive Lenia, and continuously optimized neural cellular systems can acquire task-dependent behaviour\cite{khajehabdollahi,pdnca}. Meta-learning and differentiable plasticity provide further ways of constructing systems that adapt during operation\cite{duan,wang,miconi,backpropamine}. These approaches motivate online state change, while leaving a specific question for a material learner: which observations and operations are already supplied? A controller given a reward-gradient direction and a positive movement gain may turn when a new gradient arrives without learning a new action relationship. A model given completed physical displacements can estimate an actuator mapping that a scalar-outcome learner cannot directly observe. Comparing their endpoints without specifying these inputs would confound acquired capability with supplied sensory structure.

The primary action experiment here therefore uses a deliberately restricted information boundary. Each local policy reads its completed scalar energy outcome, its transported internal state and its own exploration records. The physical transport operator carries material and memory but does not provide sensed displacement. There is no reward-gradient input, global location, field map or hidden actuator truth. Standard random perturbations can still correlate an executed action with a later local consequence. Our question is whether that restricted process supports a usable acquired direction, how its content affects a later future, and whether a fixed response to new outcomes already provides the same practical capability. These are questions about acquired action under specified physical conditions; they do not assume a novel optimizer.

Biological work underscores why retained state and retained information should also be distinguished. Habituation and learned responses in non-neural organisms, environmental traces and altered transport networks can all affect later behaviour\cite{boisseau,reid,kramar2021,vogel,boussard}. Persistent activity alone can be misleading evidence of memory\cite{reiber}. Regulatory persistence and collective adaptation create additional routes by which past conditions influence a future response\cite{acar,micheelsen,yang}. Environmental fluctuation changes the value of retaining those conditions\cite{lambert}. Rather than assigning these biological mechanisms to our artificial bodies, we use them to motivate interventions that ask what a state contains and what its removal changes.

A historical association is not sufficient for causal attribution. A body may already occupy a favourable location when its stored direction is erased. It may retain a high total association while local concentrations and motion change. Conversely, a field can be causally important without its particular spatial arrangement being necessary. Complete-state branches allow us to change an acquired quantity while matching the body's density, energy, other records and randomness at the intervention. These branches test mechanisms within an acquired history. They do not multiply the number of independent histories, and their initial agreement does not prevent subsequent physical feedback from producing different bodies.

Strong simple controls make these interventions scientifically informative. If a fixed response to worsening outcomes maintains useful movement after execution changes, continued fitting is not uniquely necessary for that ability. If a complete literal-history table solves an order-dependent task, a more elaborate representation has not established additional capacity merely by using the past. Learning can be costly\cite{mery2005}, and the relative value of experience, inherited preferences and anticipation depends on the environmental timescale\cite{dunlap2009,mitchell,tagkopoulos}. We therefore retain fixed physiological, recent-feedback and capacity-adequate history controls, while reporting their actual expenses. Equal prices, equal purchase quotas and equal information are distinct experimental conditions.

Physical continuity adds another layer. Flows and local organization can transmit or process information\cite{alimflow,alimsignal,nakajima2015}, but a quantity that travels with tissue has not thereby been shown to be a learned memory. Regeneration raises a further question: does the retained content remain useful when its original carrier is lost? Biological memory through regeneration motivates that question\cite{shomrat}, but an artificial demonstration requires both actual material replacement and an intervention on the proposed retained content. Cumulative birth, the amount of newborn tissue still present and the death of a registered old cohort answer different physical questions. None can substitute for the continuing individual's behavioural history.

We accordingly keep three core cellular model classes explicit. Kernel-driven Lenia-like generative bodies establish developmental and resource consequences under their specified rules. Conservative material bodies allow action learning and history interventions without tissue birth. An externally supplied density-regulated birth/death material CA supports real production, loss and conditional maintenance. Within the conservative class, local scalar policies, gradient-conditioned action, central histories and material-carried event memories also have different interfaces. A separate externally supported elastic-contact network tests force-equivalent controls outside these cellular models. A positive result in one system does not supply the unmeasured part of another. Allowing experience-driven behaviour to be studied beyond the original Lenia growth kernel broadens the question without merging these implementations into one organism.

The resulting evidence has three connected parts. First, scalar local experience can establish a usable direction, but its revision and later expression have finite, condition-specific limits; controls distinguish those limits from the stronger inputs supplied to earlier models. Second, interventions show how acquired direction and earlier cue order affect actual action while strong simple responses, complete history tables and fixed material-carried memory constrain claims of additional method value. Third, acquired conditional maintenance can reduce real tissue injury under equal protection purchases, yet its benefit depends on relation stability, contact timing and the physiological comparator. A representative spatial process and the resource-admission experiments connect these claims to changes in bodily organization, energy and surviving tissue.

These parts are developed as a full physical account rather than a collection of passing endpoints. A later unreinforced success does not retrospectively turn a failed reversal interval into successful relearning. Better prediction is not bodily benefit when fixed allocation achieves the same physical outcome. At the same time, useful behaviour need not increase mass, energy and every tissue cohort simultaneously: movement or preparation can have a measurable cost. The contribution of the present work is the demonstrated, bounded use of acquired content together with experiments that identify which information and physical alternatives sustain it. Whether one evolving cellular individual can acquire, revise and retain that use through its own regeneration remains an open experimental question.
